% ====================================================================
% ФАЙЛ: tasks/03_electrodynamics/coulomb_law_and_field_part1.tex
% ТЕМА: ЗАКОН КУЛОНА. НАПРЯЖЕННОСТЬ ЭЛЕКТРИЧЕСКОГО ПОЛЯ
% ПАЧКА 1: ВАРИАНТЫ 1–5
% ====================================================================

% === ЗАДАЧА 1 ===
% Тег: #ЭЛЕКТРОДИНАМИКА #КУЛОН #НАПРЯЖЕННОСТЬ #КИМ_11 #ВАРИАНТ_01
\begin{taskbasic}[code={\label{task:elem_coulomb_v1_n11}}]
\textbf{Задача 1.} Во сколько раз уменьшится модуль сил кулоновского взаимодействия двух точечных неподвижных заряженных тел, если модуль заряда каждого из тел уменьшить в $3$ раза?

\kim{11}
\end{taskbasic}
\answer{9}

% === ЗАДАЧА 2 ===
% Тег: #ЭЛЕКТРОДИНАМИКА #КУЛОН #НАПРЯЖЕННОСТЬ #КИМ_11 #ВАРИАНТ_02
\begin{taskbasic}[code={\label{task:elem_coulomb_v2_n11}}]
\textbf{Задача 2.} Во сколько раз увеличится модуль сил кулоновского взаимодействия двух точечных неподвижных заряженных тел, если расстояние между ними уменьшить в $4$ раза?

\kim{11}
\end{taskbasic}
\answer{16}

% === ЗАДАЧА 3 ===
% Тег: #ЭЛЕКТРОДИНАМИКА #КУЛОН #НАПРЯЖЕННОСТЬ #КИМ_11 #ВАРИАНТ_03
\begin{taskbasic}[code={\label{task:elem_coulomb_v3_n11}}]
\textbf{Задача 3.} В вершинах равнобедренного прямоугольного треугольника $ABC$ с катетами $a = 10\text{ см}$ расположены точечные заряды $q_A = +2\text{ нКл}$ и $q_B = -2\text{ нКл}$. Найдите модуль напряжённости электрического поля этих зарядов в точке $C$. Ответ выразите в кВ/м.

\nopagebreak\smallskip
\begin{center}
\begin{tikzpicture}[scale=1.1, every node/.style={font=\footnotesize}]
    \draw[xstep=1, ystep=1, gray!30, thin] (0,0) grid (3,3);
    \draw[-{Stealth[scale=0.8]}, black, thick] (0,0) -- (3.5,0) node[right] {$x$};
    \draw[-{Stealth[scale=0.8]}, black, thick] (0,0) -- (0,3.5) node[above] {$y$};
    
    \coordinate (A) at (1,1);
    \coordinate (B) at (3,1);
    \coordinate (C) at (1,3);
    
    \draw[thick, gray] (A) -- (B) -- (C) -- cycle;
    
    \fill[black] (A) circle (2pt) node[below left] {$A$ ($+q$)};
    \fill[black] (B) circle (2pt) node[below right] {$B$ ($-q$)};
    \fill[black] (C) circle (2pt) node[above left] {$C$};
    
    \node[left] at (0,1) {$1$};
    \node[left] at (0,3) {$3$};
    \node[below] at (1,0) {$1$};
    \node[below] at (3,0) {$3$};
    \node[below left] at (0,0) {0};
\end{tikzpicture}
\end{center}

\kim{11}
\end{taskbasic}
\answer{4,5}

% === ЗАДАЧА 4 ===
% Тег: #ЭЛЕКТРОДИНАМИКА #КУЛОН #НАПРЯЖЕННОСТЬ #КИМ_11 #ВАРИАНТ_04
\begin{taskbasic}[code={\label{task:elem_coulomb_v4_n11}}]
\textbf{Задача 4.} В двух вершинах квадрата со стороной $a = 30\text{ см}$ находятся неподвижные точечные заряды $q_1 = +4\text{ нКл}$ и $q_2 = -4\text{ нКл}$. Чему равен модуль напряжённости электрического поля этих зарядов в центре квадрата? Ответ выразите в кВ/м.

\nopagebreak\smallskip
\begin{center}
\begin{tikzpicture}[scale=1.1, every node/.style={font=\footnotesize}]
    \draw[xstep=1, ystep=1, gray!30, thin] (0,0) grid (4,4);
    \draw[-{Stealth[scale=0.8]}, black, thick] (0,0) -- (4.5,0) node[right] {$x$};
    \draw[-{Stealth[scale=0.8]}, black, thick] (0,0) -- (0,4.5) node[above] {$p$};
    
    \coordinate (1) at (1,1);
    \coordinate (2) at (3,1);
    \coordinate (3) at (3,3);
    \coordinate (4) at (1,3);
    \coordinate (O) at (2,2);
    
    \draw[thick, gray] (1) -- (2) -- (3) -- (4) -- cycle;
    \draw[dashed, gray] (1) -- (3);
    \draw[dashed, gray] (2) -- (4);
    
    \fill[black] (1) circle (2pt) node[below left] {$q_1$};
    \fill[black] (2) circle (2pt) node[below right] {$q_2$};
    \fill[black] (O) circle (2pt) node[above] {$O$};
    
    \node[below left] at (0,0) {0};
\end{tikzpicture}
\end{center}

\kim{11}
\end{taskbasic}
\answer{0,5}

% ====================================================================
% ФАЙЛ: tasks/03_electrodynamics/coulomb_law_and_field_part2.tex
% ТЕМА: ЗАКОН КУЛОНА. НАПРЯЖЕННОСТЬ ЭЛЕКТРИЧЕСКОГО ПОЛЯ
% ПАЧКА 2: ВАРИАНТЫ 6–12
% ====================================================================

% === ЗАДАЧА 1 ===
% Тег: #ЭЛЕКТРОДИНАМИКА #КУЛОН #НАПРЯЖЕННОСТЬ #КИМ_11 #ВАРИАНТ_07
\begin{taskbasic}[code={\label{task:elem_coulomb_v7_n11}}]
\textbf{Задача 1.} Во сколько раз уменьшится модуль сил кулоновского взаимодействия двух точечных неподвижных заряженных тел, если расстояние между ними увеличить в $2$ раза?

\kim{11}
\end{taskbasic}
\answer{4}

% === ЗАДАЧА 2 ===
% Тег: #ЭЛЕКТРОДИНАМИКА #КУЛОН #НАПРЯЖЕННОСТЬ #КИМ_11 #ВАРИАНТ_08
\begin{taskbasic}[code={\label{task:elem_coulomb_v8_n11}}]
\textbf{Задача 2.} Во сколько раз увеличится модуль напряжённости электрического поля точечного заряда в некоторой точке пространства, если модуль этого заряда увеличить в $2$ раза?

\kim{11}
\end{taskbasic}
\answer{2}

% === ЗАДАЧА 3 ===
% Тег: #ЭЛЕКТРОДИНАМИКА #КУЛОН #НАПРЯЖЕННОСТЬ #КИМ_14 #ВАРИАНТ_11
\begin{taskbasic}[code={\label{task:elem_coulomb_v11_n14}}]
\textbf{Задача 3.} На оси $Ox$ расположены два неподвижных точечных заряда: $q_1 = +q$ в точке $x = -a$ и $q_2 = -q$ в точке $x = +a$. Из приведённого ниже списка выберите все верные утверждения относительно электростатического поля, создаваемого этими зарядами.

\nopagebreak\smallskip
\begin{center}
\begin{tikzpicture}[scale=1.0, every node/.style={font=\footnotesize}]
    \draw[xstep=1, ystep=1, gray!30, thin] (-3,0) grid (3,1);
    \draw[-{Stealth[scale=0.8]}, black, thick] (-3.5,0) -- (3.5,0) node[right] {$x$};
    \draw[-{Stealth[scale=0.8]}, black, thick] (0,-0.5) -- (0,1.5) node[above] {$y$};
    
    \coordinate (Q1) at (-2,0);
    \coordinate (Q2) at (2,0);
    
    \fill[black] (Q1) circle (2pt) node[above] {$+q$};
    \fill[black] (Q2) circle (2pt) node[above] {$-q$};
    
    \node[below] at (-2,0) {$-a$};
    \node[below] at (2,0) {$a$};
    \node[below left] at (0,0) {0};
\end{tikzpicture}
\end{center}

1) В точке $x = 0$ вектор напряжённости суммарного электрического поля направлен вправо (вдоль оси $Ox$).\\
2) Напряжённость суммарного электрического поля в точке $x = 2a$ равна нулю.\\
3) Потенциал суммарного электрического поля в точке $x = -2a$ равен нулю.\\
4) Проекция вектора напряжённости электрического поля на ось $Ox$ в точке $x = 0$ положительна.\\
5) Модуль напряжённости электрического поля в точке $x = -2a$ больше, чем в точке $x = 0$.

\kim{14}
\end{taskbasic}
\answer{14}

% === ЗАДАЧА 4 ===
% Тег: #ЭЛЕКТРОДИНАМИКА #КУЛОН #НАПРЯЖЕННОСТЬ #КИМ_14 #ВАРИАНТ_12
\begin{taskbasic}[code={\label{task:elem_coulomb_v12_n14}}]
\textbf{Задача 4.} На оси $Ox$ расположены два одинаковых неподвижных точечных заряда: $q_1 = +q$ в точке $x = -a$ и $q_2 = +q$ в точке $x = +a$. Из приведённого ниже списка выберите все верные утверждения относительно электростатического поля, создаваемого этими зарядами.

\nopagebreak\smallskip
\begin{center}
\begin{tikzpicture}[scale=1.0, every node/.style={font=\footnotesize}]
    \draw[xstep=1, ystep=1, gray!30, thin] (-3,0) grid (3,1);
    \draw[-{Stealth[scale=0.8]}, black, thick] (-3.5,0) -- (3.5,0) node[right] {$x$};
    \draw[-{Stealth[scale=0.8]}, black, thick] (0,-0.5) -- (0,1.5) node[above] {$y$};
    
    \coordinate (Q1) at (-2,0);
    \coordinate (Q2) at (2,0);
    
    \fill[black] (Q1) circle (2pt) node[above] {$+q$};
    \fill[black] (Q2) circle (2pt) node[above] {$+q$};
    
    \node[below] at (-2,0) {$-a$};
    \node[below] at (2,0) {$a$};
    \node[below left] at (0,0) {0};
\end{tikzpicture}
\end{center}

1) В точке $x = 0$ модуль напряжённости суммарного электрического поля максимален.\\
2) В точке $x = 0$ вектор напряжённости суммарного электрического поля равен нулю.\\
3) Потенциал суммарного электрического поля в точке $x = 0$ положителен.\\
4) В точке $x = 2a$ вектор напряжённости суммарного электрического поля направлен влево.\\
5) Модуль напряжённости электрического поля в точке $x = 2a$ равен нулю.

\kim{14}
\end{taskbasic}
\answer{23}

% ====================================================================
% ФАЙЛ: tasks/03_electrodynamics/coulomb_law_and_field_part3.tex
% ТЕМА: ЗАКОН КУЛОНА. НАПРЯЖЕННОСТЬ ЭЛЕКТРИЧЕСКОГО ПОЛЯ
% ПАЧКА 3: ВАРИАНТЫ 13–22
% ====================================================================

% === ЗАДАЧА 1 ===
% Тег: #ЭЛЕКТРОДИНАМИКА #КУЛОН #НАПРЯЖЕННОСТЬ #КИМ_11 #ВАРИАНТ_15
\begin{taskbasic}[code={\label{task:elem_coulomb_v15_n11}}]
\textbf{Задача 1.} Во сколько раз уменьшится модуль сил кулоновского взаимодействия двух точечных неподвижных заряженных тел, если расстояние между ними увеличить в $3$ раза?

\kim{11}
\end{taskbasic}
\answer{9}

% === ЗАДАЧА 2 ===
% Тег: #ЭЛЕКТРОДИНАМИКА #КУЛОН #НАПРЯЖЕННОСТЬ #КИМ_11 #ВАРИАНТ_16
\begin{taskbasic}[code={\label{task:elem_coulomb_v16_n11}}]
\textbf{Задача 2.} Во сколько раз нужно уменьшить расстояние от точечного заряда до некоторой точки пространства, чтобы модуль напряжённости электрического поля в этой точке увеличился в $9$ раз?

\kim{11}
\end{taskbasic}
\answer{3}

% === ЗАДАЧА 3 ===
% Тег: #ЭЛЕКТРОДИНАМИКА #КУЛОН #НАПРЯЖЕННОСТЬ #КИМ_21 #ВАРИАНТ_21
\begin{taskbasic}[code={\label{task:elem_coulomb_v21_n21}}]
\textbf{Задача 3.} Небольшое заряженное тело массой $m$ и зарядом $q > 0$ влетает в плоский горизонтальный конденсатор с горизонтально направленной скоростью $v_0$ параллельно пластинам. Длина пластин $L$, расстояние между ними $d$. Напряжённость электрического поля внутри конденсатора равна $E$. Как изменится вертикальное смещение тела при вылете из конденсатора, если напряжённость поля $E$ увеличить в $2$ раза, а начальную скорость $v_0$ увеличить в $2$ раза? Тяготением пренебречь. Ответ обоснуйте, указав используемые физические закономерности.

\kim{21}
\end{taskbasic}
\answer{Уменьшится в 2 раза}

% === ЗАДАЧА 4 ===
% Тег: #ЭЛЕКТРОДИНАМИКА #КУЛОН #НАПРЯЖЕННОСТЬ #КИМ_21 #ВАРИАНТ_22
\begin{taskbasic}[code={\label{task:elem_coulomb_v22_n21}}]
\textbf{Задача 4.} Два положительных точечных заряда $q$ и $4q$ закреплены на расстоянии $L$ друг от друга. Где на отрезке, соединяющем эти заряды, следует поместить третий точечный заряд $-q_0$, чтобы он находился в равновесии? Укажите расстояние от заряда $q$ до заряда $-q_0$. Ответ обоснуйте, указав используемые физические закономерности.

\nopagebreak\smallskip
\begin{center}
\begin{tikzpicture}[scale=1.1, every node/.style={font=\footnotesize}]
    \draw[xstep=1, ystep=1, gray!30, thin] (0,0) grid (4,1);
    \draw[-{Stealth[scale=0.8]}, black, thick] (-0.5,0) -- (4.5,0) node[right] {$x$};
    
    \coordinate (Q1) at (0,0);
    \coordinate (Q2) at (3,0);
    
    \fill[black] (Q1) circle (2pt) node[above] {$+q$};
    \fill[black] (Q2) circle (2pt) node[above] {$+4q$};
    
    \node[below] at (0,0) {$0$};
    \node[below] at (3,0) {$L$};
\end{tikzpicture}
\end{center}

\kim{21}
\end{taskbasic}
\answer{L/3}

% ====================================================================
% ФАЙЛ: tasks/03_electrodynamics/coulomb_law_and_field_part4.tex
% ТЕМА: ЗАКОН КУЛОНА. НАПРЯЖЕННОСТЬ ЭЛЕКТРИЧЕСКОГО ПОЛЯ
% ПАЧКА 4: ВАРИАНТЫ 23–30
% ====================================================================

% === ЗАДАЧА 1 ===
% Тег: #ЭЛЕКТРОДИНАМИКА #КУЛОН #НАПРЯЖЕННОСТЬ #КИМ_11 #ВАРИАНТ_23
\begin{taskbasic}[code={\label{task:elem_coulomb_v23_n11}}]
\textbf{Задача 1.} Во сколько раз увеличится модуль сил кулоновского взаимодействия двух точечных неподвижных заряженных тел, если расстояние между ними уменьшить в $5$ раз?

\kim{11}
\end{taskbasic}
\answer{25}

% === ЗАДАЧА 2 ===
% Тег: #ЭЛЕКТРОДИНАМИКА #КУЛОН #НАПРЯЖЕННОСТЬ #КИМ_11 #ВАРИАНТ_24
\begin{taskbasic}[code={\label{task:elem_coulomb_v24_n11}}]
\textbf{Задача 2.} Во сколько раз увеличится модуль напряжённости электрического поля точечного заряда в некоторой точке пространства, если модуль этого заряда увеличить в $5$ раз?

\kim{11}
\end{taskbasic}
\answer{5}

% === ЗАДАЧА 3 ===
% Тег: #ЭЛЕКТРОДИНАМИКА #КУЛОН #НАПРЯЖЕННОСТЬ #КИМ_21 #ВАРИАНТ_25
\begin{taskbasic}[code={\label{task:elem_coulomb_v25_n21}}]
\textbf{Задача 3.} В вершинах квадрата со стороной $a$ закреплены четыре точечных заряда: трех из них равны $+q$, а один равен $-q$. Модуль напряжённости суммарного электрического поля этих зарядов в центре квадрата равен $E_0$. Чему станет равен модуль напряжённости поля в центре квадрата, если заряд $-q$ заменить на $+q$? Ответ обоснуйте, указав используемые физические закономерности.

\nopagebreak\smallskip
\begin{center}
\begin{tikzpicture}[scale=1.1, every node/.style={font=\footnotesize}]
    \draw[xstep=1, ystep=1, gray!30, thin] (0,0) grid (3,3);
    \draw[-{Stealth[scale=0.8]}, black, thick] (-0.5,0) -- (3.5,0) node[right] {$x$};
    \draw[-{Stealth[scale=0.8]}, black, thick] (0,-0.5) -- (0,3.5) node[above] {$y$};
    
    \coordinate (1) at (0.5,0.5);
    \coordinate (2) at (2.5,0.5);
    \coordinate (3) at (2.5,2.5);
    \coordinate (4) at (0.5,2.5);
    \coordinate (O) at (1.5,1.5);
    
    \draw[thick, gray] (1) -- (2) -- (3) -- (4) -- cycle;
    \draw[dashed, gray] (1) -- (3);
    \draw[dashed, gray] (2) -- (4);
    
    \fill[black] (1) circle (2pt) node[below left] {$+q$};
    \fill[black] (2) circle (2pt) node[below right] {$+q$};
    \fill[black] (3) circle (2pt) node[above right] {$-q$};
    \fill[black] (4) circle (2pt) node[above left] {$+q$};
    \fill[black] (O) circle (1.5pt) node[below] {$O$};
    
    \node[below left] at (0,0) {0};
\end{tikzpicture}
\end{center}

\kim{21}
\end{taskbasic}
\answer{0}

% === ЗАДАЧА 4 ===
% Тег: #ЭЛЕКТРОДИНАМИКА #КУЛОН #НАПРЯЖЕННОСТЬ #КИМ_25 #ВАРИАНТ_26
\begin{taskbasic}[code={\label{task:elem_coulomb_v26_n25}}]
\textbf{Задача 4.} Маленький шарик массой $m = 0{,}9\text{ г}$ с зарядом $q = 10\text{ нКл}$ подвешен на невесомой шёлковой нити между двумя вертикальными параллельными противоположно заряженными пластинами. Вектор напряжённости электрического поля между пластинами направлен горизонтально, а модуль напряжённости равен $E = 3\cdot 10^5\text{ В/м}$. Определите угол $\alpha$ отклонения нити от вертикали. Ответ выразите в градусах.

\nopagebreak\smallskip
\begin{center}
\begin{tikzpicture}[scale=1.1, every node/.style={font=\footnotesize}]
    \draw[xstep=1, ystep=1, gray!30, thin] (-1,0) grid (3,3);
    
    % Пластины
    \draw[ultra thick, black] (-0.5,0.5) -- (-0.5,2.5);
    \draw[ultra thick, black] (2.5,0.5) -- (2.5,2.5);
    \node[above] at (-0.5,2.5) {$+$};
    \node[above] at (2.5,2.5) {$-$};
    
    % Подвес и нить
    \draw[thick, black] (1,2.5) -- (1.8,1.1);
    \draw[dashed, gray] (1,2.5) -- (1,0.5);
    
    % Шарик
    \fill[brandPrimary] (1.8,1.1) circle (3pt) node[right=2pt] {$q$};
    
    % Вектор E
    \draw[-{Stealth[scale=0.8]}, thick, brandPrimary] (0,1.5) -- (0.8,1.5) node[above] {$\vec{E}$};
    
    % Угол alpha
    \draw (1,2.0) arc (270:300:0.5);
    \node at (1.15,1.8) {$\alpha$};
    
    \node[below left] at (-1,0) {0};
\end{tikzpicture}
\end{center}

\kim{25}
\end{taskbasic}
\answer{9}